\documentclass[11pt]{article}
% The preamble every project's paper shares: % The preamble every project's paper shares: % The preamble every project's paper shares: \input{../../preamble} from
% projects/<project>/paper/. Boards come from chessboard and xskak, plots from
% pgfplots; every number a paper prints comes from its generated.tex.
\usepackage[LSBC4,T1]{fontenc}
\usepackage[utf8]{inputenc}
\usepackage{lmodern}
\usepackage[margin=1in]{geometry}
\usepackage{microtype}
\usepackage{amsmath}
\usepackage{booktabs}
\usepackage{tabularx}
\usepackage{array}
\usepackage{xcolor}
\usepackage[obeyspaces,spaces]{url}
\usepackage{xskak}
\usepackage{pgfplots}
\pgfplotsset{compat=1.18}
\usepackage{caption}
\usepackage{subcaption}
\usepackage{enumitem}
\usepackage[section]{placeins}
\usepackage[hidelinks]{hyperref}

\setchessboard{boardfontencoding=LSBC4, boardfontfamily=skaknew, blackfieldcolor=black!22,
  whitefieldcolor=white, setfontcolors, showmover=false, label=true, labelfontsize=6pt,
  boardfontsize=14pt}

% A move or line in UCI, and a position in FEN, both breakable at a slash.
\newcommand{\uci}[1]{\texttt{#1}}
\newcommand{\fen}[1]{{\footnotesize\path{#1}}}
\newcolumntype{L}{>{\raggedright\arraybackslash}X}

\setlength{\parskip}{0.4em}
\setlength{\parindent}{0pt}
 from
% projects/<project>/paper/. Boards come from chessboard and xskak, plots from
% pgfplots; every number a paper prints comes from its generated.tex.
\usepackage[LSBC4,T1]{fontenc}
\usepackage[utf8]{inputenc}
\usepackage{lmodern}
\usepackage[margin=1in]{geometry}
\usepackage{microtype}
\usepackage{amsmath}
\usepackage{booktabs}
\usepackage{tabularx}
\usepackage{array}
\usepackage{xcolor}
\usepackage[obeyspaces,spaces]{url}
\usepackage{xskak}
\usepackage{pgfplots}
\pgfplotsset{compat=1.18}
\usepackage{caption}
\usepackage{subcaption}
\usepackage{enumitem}
\usepackage[section]{placeins}
\usepackage[hidelinks]{hyperref}

\setchessboard{boardfontencoding=LSBC4, boardfontfamily=skaknew, blackfieldcolor=black!22,
  whitefieldcolor=white, setfontcolors, showmover=false, label=true, labelfontsize=6pt,
  boardfontsize=14pt}

% A move or line in UCI, and a position in FEN, both breakable at a slash.
\newcommand{\uci}[1]{\texttt{#1}}
\newcommand{\fen}[1]{{\footnotesize\path{#1}}}
\newcolumntype{L}{>{\raggedright\arraybackslash}X}

\setlength{\parskip}{0.4em}
\setlength{\parindent}{0pt}
 from
% projects/<project>/paper/. Boards come from chessboard and xskak, plots from
% pgfplots; every number a paper prints comes from its generated.tex.
\usepackage[LSBC4,T1]{fontenc}
\usepackage[utf8]{inputenc}
\usepackage{lmodern}
\usepackage[margin=1in]{geometry}
\usepackage{microtype}
\usepackage{amsmath}
\usepackage{booktabs}
\usepackage{tabularx}
\usepackage{array}
\usepackage{xcolor}
\usepackage[obeyspaces,spaces]{url}
\usepackage{xskak}
\usepackage{pgfplots}
\pgfplotsset{compat=1.18}
\usepackage{caption}
\usepackage{subcaption}
\usepackage{enumitem}
\usepackage[section]{placeins}
\usepackage[hidelinks]{hyperref}

\setchessboard{boardfontencoding=LSBC4, boardfontfamily=skaknew, blackfieldcolor=black!22,
  whitefieldcolor=white, setfontcolors, showmover=false, label=true, labelfontsize=6pt,
  boardfontsize=14pt}

% A move or line in UCI, and a position in FEN, both breakable at a slash.
\newcommand{\uci}[1]{\texttt{#1}}
\newcommand{\fen}[1]{{\footnotesize\path{#1}}}
\newcolumntype{L}{>{\raggedright\arraybackslash}X}

\setlength{\parskip}{0.4em}
\setlength{\parindent}{0pt}

% Written by projects/explain-top-move/paper/generate.mjs; do not edit by hand.

\newcommand{\MomentsTest}{152}
\newcommand{\MomentsTrain}{177}
\newcommand{\Moments}{329}
\newcommand{\GamesTest}{30}
\newcommand{\GamesTrain}{30}
\newcommand{\Games}{60}
\newcommand{\HandcheckMoments}{19}
\newcommand{\PlayerRatingMin}{948}
\newcommand{\PlayerRatingMax}{1167}
\newcommand{\HandcheckAgree}{19}
\newcommand{\BaseTestFaithful}{90.1}
\newcommand{\BaseTestFaithfulCI}{90.1 [84.2, 95.9]}
\newcommand{\BaseTestConcrete}{75.0}
\newcommand{\BaseTestConcreteCI}{75.0 [65.3, 83.6]}
\newcommand{\BaseTestNamedDeep}{94.1}
\newcommand{\BaseTestNamedDeepCI}{94.1 [89.7, 97.9]}
\newcommand{\BaseTestFaithfulAll}{89.5}
\newcommand{\BaseTestFaithfulAllCI}{89.5 [83.3, 95.2]}
\newcommand{\KeptTestFaithful}{100.0}
\newcommand{\KeptTestFaithfulCI}{100.0 [100.0, 100.0]}
\newcommand{\KeptTestConcrete}{82.9}
\newcommand{\KeptTestConcreteCI}{82.9 [72.0, 91.4]}
\newcommand{\KeptTestNamedDeep}{94.1}
\newcommand{\KeptTestNamedDeepCI}{94.1 [89.7, 97.9]}
\newcommand{\KeptTestFaithfulAll}{99.3}
\newcommand{\KeptTestFaithfulAllCI}{99.3 [97.8, 100.0]}
\newcommand{\KeptTestFaithfulVsBase}{+9.9 [4.0, 15.9]}
\newcommand{\KeptTestConcreteVsBase}{+7.9 [3.9, 12.1]}
\newcommand{\KeptTestTextsChanged}{70}
\newcommand{\BaseTrainFaithful}{87.6}
\newcommand{\BaseTrainFaithfulCI}{87.6 [83.0, 91.7]}
\newcommand{\BaseTrainConcrete}{78.0}
\newcommand{\BaseTrainConcreteCI}{78.0 [70.7, 84.2]}
\newcommand{\BaseTrainNamedDeep}{92.1}
\newcommand{\BaseTrainNamedDeepCI}{92.1 [86.8, 97.1]}
\newcommand{\BaseTrainFaithfulAll}{87.6}
\newcommand{\BaseTrainFaithfulAllCI}{87.6 [82.9, 91.8]}
\newcommand{\KeptTrainFaithful}{99.4}
\newcommand{\KeptTrainFaithfulCI}{99.4 [98.1, 100.0]}
\newcommand{\KeptTrainConcrete}{84.7}
\newcommand{\KeptTrainConcreteCI}{84.7 [78.9, 89.8]}
\newcommand{\KeptTrainNamedDeep}{92.1}
\newcommand{\KeptTrainNamedDeepCI}{92.1 [86.8, 97.1]}
\newcommand{\KeptTrainFaithfulAll}{99.4}
\newcommand{\KeptTrainFaithfulAllCI}{99.4 [98.2, 100.0]}
\newcommand{\KeptTrainFaithfulVsBase}{+11.9 [7.9, 16.4]}
\newcommand{\KeptTrainConcreteVsBase}{+6.8 [3.2, 10.6]}
\newcommand{\KeptTrainTextsChanged}{87}
\newcommand{\KeptVariant}{v8}
\newcommand{\RoundCount}{nine}
\newcommand{\KeptRoundCount}{seven}
\newcommand{\RoundsTable}{%
\begin{tabularx}{\textwidth}{l L rrrr r r}
\toprule
 & & \multicolumn{2}{c}{test (\%)} & \multicolumn{2}{c}{train (\%)} & texts & judge \\
\cmidrule(lr){3-4}\cmidrule(lr){5-6}
variant & change & faithful & concrete & faithful & concrete & changed & train / test \\
\midrule
baseline & the app before the rounds & 90.1 & 75.0 & 87.6 & 78.0 &  &  \\
v1 & Say each line's material as it is in the lesson's difference & 96.1 & 82.2 & 98.3 & 84.7 & 20, 34 &  \\
v2 & Name only the pawns the seventh-rank rook hits & 100.0 & 82.2 & 99.4 & 84.7 & 25, 35 &  \\
v3 & Say what the better move defends & 100.0 & 82.2 & 99.4 & 84.7 & 31, 43 & 22--0 / 18--2 \\
v4 & Answer the capture a learner asks about first & 100.0 & 82.9 & 99.4 & 84.7 & 34, 45 & 8--0 / 6--0 \\
v5 & Say where a long mate ends (\emph{not kept}) & 100.0 & 82.9 & 99.4 & 84.7 & 45, 54 & 6--6 / 14--2 \\
v6 & Say where a long mate ends, when the end is in reach & 100.0 & 82.9 & 99.4 & 84.7 & 44, 51 & 4--0 / 10--0 \\
v7 & Say what a capture the better line declines would cost (\emph{replaced by v7b}) & 98.7 & 82.9 & 97.2 & 84.7 & 65, 77 & 48--3 / 40--1 \\
v7b & Say what a capture wins, not the balance it leaves & 100.0 & 82.9 & 99.4 & 84.7 & 65, 77 & 51--0 / 41--0 \\
v8 & Say when the played move is right but too early & 100.0 & 82.9 & 99.4 & 84.7 & 70, 87 & 50--0 / 20--0 \\
v9 & Answer the natural capture, and skip what a pawn plainly guards (\emph{not kept}) & 100.0 & 82.2 & 98.9 & 85.9 & 142, 166 & 1--35 / 0--12 \\
\bottomrule
\end{tabularx}
}
\newcommand{\JudgeTable}{%
\begin{tabular}{l l rr l}
\toprule
round & judged & train new--old (same) & test new--old (same) & the rule says \\
\midrule
3 & round 3 & 22--0 (0) & 18--2 (0) & keep \\
4 & round 4 & 8--0 (0) & 6--0 (0) & keep \\
5 & round 5 & 6--6 (10) & 14--2 (12) & do not keep \\
6 & round 5 (17 moments) & 4--0 (8) & 10--0 (12) & keep \\
7 & round 7 & 48--3 (15) & 40--1 (9) & keep \\
7b & round 7b & 20--0 (0) & 12--0 (0) & keep \\
7 and 7b & round 7 then round 7b & 51--0 (15) & 41--0 (9) & keep \\
8 & round 8 & 50--0 (0) & 20--0 (0) & keep \\
9 & round 9 & 1--35 (0) & 0--12 (0) & do not keep \\
\bottomrule
\end{tabular}
}
\newcommand{\RoundFiveTrainNew}{6}
\newcommand{\RoundFiveTrainOld}{6}
\newcommand{\RoundFiveTestNew}{14}
\newcommand{\RoundFiveTestOld}{2}
\newcommand{\RoundFiveReason}{train 4 new vs 6 old (test 10 vs 2): the picture hurt mates of 8 to 13 moves}
\newcommand{\ExampleARound}{3}
\newcommand{\ExampleACaption}{Round 3: the judge preferred the new text in both runs. The move played is the red arrow, the engine's best move the green one.}
\newcommand{\ExampleA}{% Written by projects/explain-top-move/paper/generate.mjs; do not edit by hand.
\begin{minipage}[t]{0.4\textwidth}\vspace{0pt}
\chessboard[setfen={1rr3k1/5pp1/p1pbpq1p/P2p3b/Q2P2P1/2P2N1P/1P1N1P2/R3R1K1 b - - 0 18}, boardfontsize=13pt, pgfstyle=straightmove, arrow=to, linewidth=0.07em, color=red!75!black, markmoves={b8-b2}, pgfstyle=straightmove, arrow=to, linewidth=0.07em, color=green!50!black, markmoves={h5-g6}]
\end{minipage}\hfill
\begin{minipage}[t]{0.56\textwidth}\vspace{0pt}\small
\textbf{Played:} 18\ldots{}\symrook{}xb2, Black to move; the app's verdict: Mistake.\par
\textbf{Engine:} +1.8 with the best move, +0.1 after yours.\par
\textbf{Best line:} 18\ldots{}\symbishop{}g6 19.~b4 h5 20.~\symknight{}e5 hxg4 21.~\symknight{}xg4\par
\textbf{UCI:} \uci{b8b2}; best line \uci{h5g6 b2b4 h6h5 f3e5 h5g4 e5g4}\par
\textbf{FEN:} \fen{1rr3k1/5pp1/p1pbpq1p/P2p3b/Q2P2P1/2P2N1P/1P1N1P2/R3R1K1 b - - 0 18}
\end{minipage}

\smallskip\small
\begin{tabularx}{\textwidth}{>{\raggedright\arraybackslash}p{0.17\textwidth} L L}
\toprule
section & before & after \\
\midrule
On the card & Rxb2 allows gxh5, and it wins the bishop outright. You lose a pawn in the engine's line. Bg6 avoids it. Mistake: it cost 15\% of your winning chances. & Rxb2 allows gxh5, and it wins the bishop outright. You lose a pawn in the engine's line. Bg6 moves your bishop on h5 away, where it had no defender. Mistake: it cost 15\% of your winning chances. \\
What Bg6 does instead & 18...Bg6 19.b4 h5 20.Ne5 hxg4 21.Nxg4 Qh4 22.Kg2 & it moves your bishop on h5 away, where it had no defender \textperiodcentered{} 18...Bg6 19.b4 h5 20.Ne5 hxg4 21.Nxg4 Qh4 22.Kg2 \\
The difference & After Rxb2, they win a pawn in the engine's line. Bg6 keeps it. & After Rxb2, they win a pawn in the engine's line. Bg6 keeps it: it moves your bishop on h5 away, where it had no defender. \\
Retry hint & \emph{(not shown)} & Hint: it moves your bishop on h5 away, where it had no defender. The move starts on the lit square. \\
Retry answer & \emph{(not shown)} & The move was Bg6. it moves your bishop on h5 away, where it had no defender \\
\bottomrule
\end{tabularx}

\textbf{The judge (run 1, where the new text was Y):} ``Y says Bg6 rescues the attacked, undefended bishop on h5 that gxh5 wins, where X only says Bg6 'avoids it'.''
}
\newcommand{\ExampleBRound}{4}
\newcommand{\ExampleBCaption}{Round 4: the judge preferred the new text in both runs. The move played is the red arrow, the engine's best move the green one.}
\newcommand{\ExampleB}{% Written by projects/explain-top-move/paper/generate.mjs; do not edit by hand.
\begin{minipage}[t]{0.4\textwidth}\vspace{0pt}
\chessboard[setfen={4r1k1/p2r2p1/5p1p/1R6/8/1P6/P2n1BPP/5RK1 w - - 1 29}, boardfontsize=13pt, pgfstyle=straightmove, arrow=to, linewidth=0.07em, color=red!75!black, markmoves={f1-d1}, pgfstyle=straightmove, arrow=to, linewidth=0.07em, color=green!50!black, markmoves={f1-e1}]
\end{minipage}\hfill
\begin{minipage}[t]{0.56\textwidth}\vspace{0pt}\small
\textbf{Played:} 29.~\symrook{}d1, White to move; the app's verdict: Blunder.\par
\textbf{Engine:} +0.0 with the best move, $-$4.7 after yours.\par
\textbf{Best line:} 29.~\symrook{}e1 \symknight{}e4 30.~\symrook{}b4 f5 31.~\symrook{}a4 \symrook{}ee7\par
\textbf{UCI:} \uci{f1d1}; best line \uci{f1e1 d2e4 b5b4 f6f5 b4a4 e8e7}\par
\textbf{FEN:} \fen{4r1k1/p2r2p1/5p1p/1R6/8/1P6/P2n1BPP/5RK1 w - - 1 29}
\end{minipage}

\smallskip\small
\begin{tabularx}{\textwidth}{>{\raggedright\arraybackslash}p{0.17\textwidth} L L}
\toprule
section & before & after \\
\midrule
On the card & Rd1 allows Nf3+, and it attacks the pawn on h2. You lose a piece in the engine's line. Re1 attacks the rook on e8 (undefended). Blunder: it cost 35\% of your winning chances. & Rd1 allows Nf3+, and if you take it with gxf3, Rxd1+ takes your rook on d1. You lose a piece in the engine's line. Re1 attacks the rook on e8 (undefended). Blunder: it cost 35\% of your winning chances. \\
Takeaway & Before you move, ask what their best reply attacks. Nf3+ was the answer to Rd1, and it attacks the pawn on h2. & Before you move, ask what their best reply attacks. Nf3+ was the answer to Rd1, and it if you take it with gxf3, Rxd1+ takes your rook on d1. \\
What Rd1 allowed: Nf3+ & it attacks the pawn on h2; it comes with check \textperiodcentered{} 29.Rd1 Nf3+ 30.gxf3 Rxd1+ 31.Kg2 Rd2 32.Kg3 & if you take it with gxf3, Rxd1+ takes your rook on d1; it attacks the pawn on h2; it comes with check \textperiodcentered{} 29.Rd1 Nf3+ 30.gxf3 Rxd1+ 31.Kg2 Rd2 32.Kg3 \\
\bottomrule
\end{tabularx}

\textbf{The judge (run 1, where the new text was X):} ``X names the real tactic, that Nf3+ comes with check and uncovers the d7 rook so Rxd1+ wins your rook on d1, while Y only says the knight attacks h2, which does not explain the lost material.''
}
\newcommand{\ExampleCRound}{7}
\newcommand{\ExampleCCaption}{Round 7: the judge preferred the new text in both runs. The move played is the red arrow, the engine's best move the green one.}
\newcommand{\ExampleC}{% Written by projects/explain-top-move/paper/generate.mjs; do not edit by hand.
\begin{minipage}[t]{0.4\textwidth}\vspace{0pt}
\chessboard[setfen={rn1q2k1/ppp3p1/3b1p1p/3p1R1Q/3P3B/2P1r3/PP1N2PP/R5K1 w - - 0 17}, boardfontsize=13pt, pgfstyle=straightmove, arrow=to, linewidth=0.07em, color=red!75!black, markmoves={a1-f1}, pgfstyle=straightmove, arrow=to, linewidth=0.07em, color=green!50!black, markmoves={f5-f6}]
\end{minipage}\hfill
\begin{minipage}[t]{0.56\textwidth}\vspace{0pt}\small
\textbf{Played:} 17.~\symrook{}af1, White to move; the app's verdict: Blunder.\par
\textbf{Engine:} +2.5 with the best move, 0.0 after yours.\par
\textbf{Best line:} 17.~\symrook{}xf6 \symqueen{}e8 18.~\symqueen{}xd5+ \symking{}h8 19.~\symrook{}f7 \symknight{}c6\par
\textbf{UCI:} \uci{a1f1}; best line \uci{f5f6 d8e8 h5d5 g8h8 f6f7 b8c6}\par
\textbf{FEN:} \fen{rn1q2k1/ppp3p1/3b1p1p/3p1R1Q/3P3B/2P1r3/PP1N2PP/R5K1 w - - 0 17}
\end{minipage}

\smallskip\small
\begin{tabularx}{\textwidth}{>{\raggedright\arraybackslash}p{0.17\textwidth} L L}
\toprule
section & before & after \\
\midrule
What Rxf6 does instead & it attacks the bishop on d6 and the pawn on h6 \textperiodcentered{} 17.Rxf6 Qe8 18.Qxd5+ Kh8 19.Rf7 Nc6 20.Raf1 Ne7 & it attacks the bishop on d6 and the pawn on h6; if they take it with gxf6, Qg6+ Kf8 Qxh6+ Kf7, then Qxe3 takes their rook on e3, and you are two pawns' worth of material up; if they take it with Qxf6, Bxf6 wins their queen for your rook \textperiodcentered{} 17.Rxf6 Qe8 18.Qxd5+ Kh8 19.Rf7 Nc6 20.Raf1 Ne7 \\
Retry answer & The move was Rxf6 (the green arrow). It attacks the bishop on d6 and the pawn on h6. & The move was Rxf6 (the green arrow). It attacks the bishop on d6 and the pawn on h6. If they take it with gxf6, Qg6+ Kf8 Qxh6+ Kf7, then Qxe3 takes their rook on e3, and you are two pawns' worth of material up. If they take it with Qxf6, Bxf6 wins their queen for your rook. \\
\bottomrule
\end{tabularx}

\textbf{The judge (run 1, where the new text was X):} ``X shows why the rook sacrifice works (gxf6 Qg6+ Kf8 Qxh6+ then Qxe3 wins the rook back, and Qxf6 Bxf6 wins the queen), while Y leaves the player wondering why Rxf6 does not just lose a rook to gxf6.''
}
\newcommand{\ExampleDRound}{8}
\newcommand{\ExampleDCaption}{Round 8: the judge preferred the new text in both runs. The move played is the red arrow, the engine's best move the green one.}
\newcommand{\ExampleD}{% Written by projects/explain-top-move/paper/generate.mjs; do not edit by hand.
\begin{minipage}[t]{0.4\textwidth}\vspace{0pt}
\chessboard[setfen={rn1q2k1/ppp1r1p1/3b1p1p/3pNp1Q/3Pn2B/2P1P3/PP3PPP/RN3RK1 w - - 0 13}, boardfontsize=13pt, pgfstyle=straightmove, arrow=to, linewidth=0.07em, color=red!75!black, markmoves={f2-f3}, pgfstyle=straightmove, arrow=to, linewidth=0.07em, color=green!50!black, markmoves={e5-g6}]
\end{minipage}\hfill
\begin{minipage}[t]{0.56\textwidth}\vspace{0pt}\small
\textbf{Played:} 13.~f3, White to move; the app's verdict: Inaccuracy.\par
\textbf{Engine:} +1.5 with the best move, +0.5 after yours.\par
\textbf{Best line:} 13.~\symknight{}g6 \symrook{}e8 14.~\symqueen{}xf5 c6 15.~f3 \symknight{}g5\par
\textbf{UCI:} \uci{f2f3}; best line \uci{e5g6 e7e8 h5f5 c7c6 f2f3 e4g5}\par
\textbf{FEN:} \fen{rn1q2k1/ppp1r1p1/3b1p1p/3pNp1Q/3Pn2B/2P1P3/PP3PPP/RN3RK1 w - - 0 13}
\end{minipage}

\smallskip\small
\begin{tabularx}{\textwidth}{>{\raggedright\arraybackslash}p{0.17\textwidth} L L}
\toprule
section & before & after \\
\midrule
Takeaway & Look for the more active move first: Ng6 attacks the rook on e7. f3 can often still be played a move later. & Look for the more active move first: Ng6 attacks the rook on e7. f3 still comes, two moves later, in the engine's line. \\
The difference & You win a pawn in the next 5 moves either way, so the difference is positional: Ng6 attacks the rook on e7. & You win a pawn in the next 5 moves either way, so the difference is positional: Ng6 attacks the rook on e7. The order matters: f3 comes in the engine's line too, two moves later (Ng6 Re8 Qxf5 c6 f3). \\
\bottomrule
\end{tabularx}

\textbf{The judge (run 1, where the new text was X):} ``X says correctly that f3 still comes two moves later in the engine's line (Ng6 Re8 Qxf5 c6 f3), which shows the problem is move order; Y's generic 'can often still be played a move later' is vaguer and gets the count wrong.''
}
\newcommand{\ExampleERound}{9}
\newcommand{\ExampleECaption}{Round 9: the judge preferred the old text in both runs. The move played is the red arrow, the engine's best move the green one.}
\newcommand{\ExampleE}{% Written by projects/explain-top-move/paper/generate.mjs; do not edit by hand.
\begin{minipage}[t]{0.4\textwidth}\vspace{0pt}
\chessboard[setfen={r3r1k1/p5p1/3R1p1p/np6/3P3B/2P5/PP4PP/5RK1 w - - 1 23}, boardfontsize=13pt, pgfstyle=straightmove, arrow=to, linewidth=0.07em, color=red!75!black, markmoves={d6-d5}, pgfstyle=straightmove, arrow=to, linewidth=0.07em, color=green!50!black, markmoves={b2-b3}]
\end{minipage}\hfill
\begin{minipage}[t]{0.56\textwidth}\vspace{0pt}\small
\textbf{Played:} 23.~\symrook{}d5, White to move; the app's verdict: Inaccuracy.\par
\textbf{Engine:} +2.1 with the best move, +1.0 after yours.\par
\textbf{Best line:} 23.~b3 \symrook{}e2 24.~\symrook{}a6 \symknight{}b7 25.~\symrook{}e1 \symrook{}c2\par
\textbf{UCI:} \uci{d6d5}; best line \uci{b2b3 e8e2 d6a6 a5b7 f1e1 e2c2}\par
\textbf{FEN:} \fen{r3r1k1/p5p1/3R1p1p/np6/3P3B/2P5/PP4PP/5RK1 w - - 1 23}
\end{minipage}

\smallskip\small
\begin{tabularx}{\textwidth}{>{\raggedright\arraybackslash}p{0.17\textwidth} L L}
\toprule
section & before & after \\
\midrule
What b3 does instead & if they take it with Nxb3, axb3 wins their knight for your pawn \textperiodcentered{} 23.b3 Re2 24.Ra6 Nb7 25.Re1 Rc2 26.Re7 Nd8 & 23.b3 Re2 24.Ra6 Nb7 25.Re1 Rc2 26.Re7 Nd8 \\
Retry answer & The move was b3 (the green arrow). If they take it with Nxb3, axb3 wins their knight for your pawn. & The move was b3 (the green arrow). \\
\bottomrule
\end{tabularx}

\textbf{The judge (run 1, where the new text was X):} ``Y adds the true safety check that b3 does not hang to the a5 knight, since Nxb3 axb3 wins the knight.''
}
\newcommand{\ExampleCount}{5}
\newcommand{\ReproLinesMatched}{27}
\newcommand{\ReproLines}{27}
\newcommand{\AppCommit}{\texttt{d8b7712}}
\newcommand{\RecordsCommit}{\texttt{b3b827d}}
\newcommand{\RerunProbes}{88}
\newcommand{\RerunEvidence}{review-evidence/5}
\newcommand{\RerunTestMoments}{152}
\newcommand{\RerunTestShared}{146}
\newcommand{\RerunTestSameGrades}{130}
\newcommand{\RerunTestFaithful}{100.0}
\newcommand{\RerunTestConcrete}{82.9}
\newcommand{\RerunTrainMoments}{177}
\newcommand{\RerunTrainShared}{171}
\newcommand{\RerunTrainSameGrades}{160}
\newcommand{\RerunTrainFaithful}{99.4}
\newcommand{\RerunTrainConcrete}{83.6}
\newcommand{\ReproTable}{%
\begin{tabular}{ll lll}
\toprule
split & grade & recorded baseline & recorded v8 & pinned app, rerun \\
\midrule
test & faithful & 90.1 [84.2, 95.9] & 100.0 [100.0, 100.0] & 100.0 [100.0, 100.0] \\
test & concrete & 75.0 [65.3, 83.6] & 82.9 [72.0, 91.4] & 82.9 [73.9, 90.5] \\
test & named\_deep & 94.1 [89.7, 97.9] & 94.1 [89.7, 97.9] & 94.1 [89.7, 97.8] \\
test & faithfulAll & 89.5 [83.3, 95.2] & 99.3 [97.8, 100.0] & 99.3 [97.7, 100.0] \\
train & faithful & 87.6 [83.0, 91.7] & 99.4 [98.1, 100.0] & 99.4 [98.1, 100.0] \\
train & concrete & 78.0 [70.7, 84.2] & 84.7 [78.9, 89.8] & 83.6 [77.0, 89.4] \\
train & named\_deep & 92.1 [86.8, 97.1] & 92.1 [86.8, 97.1] & 92.1 [86.7, 97.2] \\
train & faithfulAll & 87.6 [82.9, 91.8] & 99.4 [98.2, 100.0] & 99.4 [98.1, 100.0] \\
\bottomrule
\end{tabular}
}


\title{Explaining the better move\\[0.3em]
\large Nine rounds of hill-climbing a game review's lessons, graded on the board and by a blind judge}
\author{Chess Coach research lab}
\date{}

\begin{document}
\maketitle

\begin{abstract}
At each key moment of a game review, the coaching app explains what the move played allowed,
what the engine's better move does instead, and the difference between them. We measured
those explanations on the \Moments{} key moments of \Games{} of the player's games, with three
grades checked against the board, and changed the code that writes them one idea per round,
keeping a round only when no grade fell on either half of the games. Where a change wrote
something the grades cannot check, a judge model compared the old and new texts blind, twice,
with the order reversed. Over \RoundCount{} rounds (\KeptRoundCount{} kept), the share of held-out moments whose
texts make no false checkable claim rose from \BaseTestFaithful\% to \KeptTestFaithful\%, and
the share that names a concrete consequence of the better move from \BaseTestConcrete\% to
\KeptTestConcrete\%. The judge preferred every kept round's texts in nearly every comparison
it decided. The kept code is the app's. The lab recomputes every published number from the
records, and the app rerun on fresh engine evidence scores \RerunTestFaithful\% and
\RerunTestConcrete\% on the same two measures.
\end{abstract}

\section{What the player reads}

When a game review stops at a key moment (a mistake, a blunder, a missed win), the app shows a
card with a headline and a verdict, a lesson panel (what the move played allowed, what the
better move does instead, the difference between the two lines, a takeaway), a Why sheet with
the evidence behind the verdict, and a retry: the player may try the position again, with a
hint and then the answer. Every one of those texts is written by heuristics in the app's review
and knowledge services from the engine's evidence. No language model writes them, so every
claim they make can be traced to a template, and most can be checked on the board.

The player is an adult club player, rated \PlayerRatingMin--\PlayerRatingMax{} on chess.com across the
\Games{} games studied here. Two things matter: that the texts are true, and that they say
\emph{why} the better move is better, in terms the player can use at the board, not only by how
much.

\section{Data and measures}

\paragraph{Moments.} The \Moments{} key moments the review stops at in \Games{} of the player's
games, from the app's phone-profile engine evidence for each game. The games are split
\GamesTrain{} for train and \GamesTest{} for test, the same split as the accuracy project, so no
game is train in one and test in the other: \MomentsTrain{} train moments and \MomentsTest{}
test moments. A round was designed on train moments only.

\paragraph{Grades.} Each moment's texts get three grades, all checked against the board with
chess.js and against the stored engine lines (\texttt{grade.mjs}):

\begin{itemize}[nosep]
\item \textbf{faithful}: 1 unless a checkable claim in the visible text is false: an illegal
move, a line that does not replay, a piece said to attack what it does not, material said to
be won that the line does not win, a mate that is not there, or an ``if they take it with''
answer that is not the engine's line or whose claim does not hold at its end.
\item \textbf{concrete}: 1 when the better move's explanation names a consequence (material, a
mate, a named threat, a prevented reply) and the text says what the move played allowed or
missed.
\item \textbf{named\_deep}: 1 when every move the text names as best rests on a search at the
app's best budget. It is a property of the evidence, not of the words, and is held, not
targeted.
\end{itemize}

A fourth grade, \textbf{faithfulAll}, is faithful over the lesson the Why sheet also shows on a
moment whose card shows only a comment. A clause no template recognises is reported as
unparsed, never graded true. The grader was checked by hand on \HandcheckMoments{} moments
before any round; the hand verdicts and the baseline's grades agree on all three grades for
\HandcheckAgree{} of them.

Every rate below is a share of key moments, with a 95\% interval from a bootstrap over games
(a game's moments share its evidence, so the game is the unit). A variant's difference from the
baseline is paired on the same moments.

\paragraph{The judge.} The grades cannot see whether a new clause \emph{helps}. For a round that
writes something new, every moment whose texts change is shown to a judge model as two versions
(X and Y, by a seeded coin), each laid out as the app lays it out, with the position, the move,
the verdict, the engine's evaluation and best line. The judge answers which version would help
the player more (understand why the better move is better, and know what to look for next
time), or that they are the same, with one sentence of reason. It answers twice, with X and Y
swapped the second time. The rule, fixed before any run: keep the round if, pooling both runs,
the new texts are preferred in more than half of the decided comparisons on train \emph{and} on
test.

\section{The rounds}

Table~\ref{tab:rounds} lists every variant. Rounds 1 and 2 fix false claims the baseline's
templates made, found by the grader over all moments before the split was used to design
anything. Rounds 3 to 9 add something the player asked for on the app's Coach Calibration page or
that the judge's notes named, so the judge decides them.

\begin{table}[!htbp]
\centering\footnotesize
\RoundsTable
\caption{Every variant. Rates are shares of key moments; \emph{texts changed} counts the test and
train moments whose visible texts differ from the baseline's; \emph{judge} is the round's pooled
preferences, new--old, on train and test (Table~\ref{tab:judge}).}
\label{tab:rounds}
\end{table}

\begin{description}[style=nextline, nosep]
\item[Round 1: each line's material as it is.] The lesson's fallback sentence said ``neither
line wins material in the next 5 moves'' whenever no earlier branch fired, including when both
lines won the same material. It now says each line's material.
\item[Round 2: the seventh-rank rook.] The insight named pawns the rook could not see past
another pawn. It now names only the pawns the rook hits.
\item[Round 3: what the better move defends.] A piece the reply takes, guarded or moved away by
the better move (Figure~\ref{fig:exa}).
\item[Round 4: the capture a learner asks about first.] Where the better move or the reply lands
on a square it can be taken on and the engine's line takes it, the lesson says what follows
(Figure~\ref{fig:exb}).
\item[Rounds 5 and 6: where a long mate ends.] A mate is told with where it ends and what holds
it there. Round 5 did so for every mate of three moves or more and failed its rule on train; round
6 keeps it for mates of three to seven moves with no promotion on the way.
\item[Round 7: a capture the better line declines.] A search of each capture the line declines
(\texttt{probes.mjs}, at the routine's budget), and the lesson says what it would cost: a
recapture as the trade it is, or checks that win a piece, move by move (Figure~\ref{fig:exc}).
Round 7b fixes the material clause the judge found false in a few moments (a gain read as a
balance).
\item[Round 8: the right move, too early.] Where the move played is in the engine's best line at
the mover's next or second turn, the difference and the takeaway say so, with the line
(Figure~\ref{fig:exd}).
\item[Round 9 (not kept): fewer capture answers.] Skipping an answer when a pawn plainly guards
the square, and stopping at the first capture with no answer. The judge wanted every answer back
(Figure~\ref{fig:exe}).
\end{description}

\begin{table}[!htbp]
\centering\footnotesize
\JudgeTable
\caption{The judge rounds, pooled over both runs: new--old preferred, with the same-answers in
parentheses. Round 6 made no run of its own: it reads round 5's judgments on the moments whose
texts round 6 changes.}
\label{tab:judge}
\end{table}

\section{Examples}

Each example is the first moment of its round, in the round's own order, on which both runs of
the judge agreed. The tables show only the sections whose text the round changed.

\begin{figure}[!htbp]\ExampleA\caption{\ExampleACaption}\label{fig:exa}\end{figure}
\begin{figure}[!htbp]\ExampleB\caption{\ExampleBCaption}\label{fig:exb}\end{figure}
\begin{figure}[!htbp]\ExampleC\caption{\ExampleCCaption}\label{fig:exc}\end{figure}
\begin{figure}[!htbp]\ExampleD\caption{\ExampleDCaption}\label{fig:exd}\end{figure}
\begin{figure}[!htbp]\ExampleE\caption{\ExampleECaption}\label{fig:exe}\end{figure}

\section{Results}

From the baseline to the kept variant (\KeptVariant), on the held-out games:

\begin{itemize}[nosep]
\item faithful: \BaseTestFaithfulCI\% to \KeptTestFaithfulCI\%, a paired change of
\KeptTestFaithfulVsBase{} points;
\item concrete: \BaseTestConcreteCI\% to \KeptTestConcreteCI\%, a paired change of
\KeptTestConcreteVsBase{} points;
\item named\_deep: \BaseTestNamedDeepCI\% throughout, as expected of a property of the evidence;
\item faithfulAll: \BaseTestFaithfulAllCI\% to \KeptTestFaithfulAllCI\%.
\end{itemize}

On train the same changes read \BaseTrainFaithful\% to \KeptTrainFaithful\% (faithful,
\KeptTrainFaithfulVsBase) and \BaseTrainConcrete\% to \KeptTrainConcrete\% (concrete,
\KeptTrainConcreteVsBase). The kept variant changes the visible texts of \KeptTestTextsChanged{}
test and \KeptTrainTextsChanged{} train moments. Every kept round past the second was preferred
by the judge on both halves; round 5 was not, and round 9 was rejected nearly unanimously.

\section{Reproduction in the lab}

\paragraph{Recomputed.} The lab recomputes every variant's aggregates from the per-moment grades
the records keep (\texttt{summarize.mjs}): \ReproLinesMatched{} of the \ReproLines{} lines of the
published aggregates match exactly. Every judge round's tally is recomputed from the judge's
answers and the order keys (\texttt{judge/pairwise.mjs}), and every decision follows from the
rule as recorded. One recorded reason misstates its counts: round 5's reads ``\RoundFiveReason'',
while the recorded answers give \RoundFiveTrainNew--\RoundFiveTrainOld{} on train and
\RoundFiveTestNew--\RoundFiveTestOld{} on test. The new counts it gives are those of the moments
round 6 later kept; the decision is the same, since a tie on train is not a majority.

\paragraph{Rerun.} The lab rebuilt the engine evidence for the same games on the app as it now runs
(\AppCommit, evidence \RerunEvidence), redid the capture searches on it (\RerunProbes{} kept), and
graded the app's texts. The rebuilt evidence is not the hill-climb's: the review pipeline changed
after the rounds, so engine lines differ in detail and some key moments move (on test,
\RerunTestShared{} of \RerunTestMoments{} moments are the same moments; on train
\RerunTrainShared{} of \RerunTrainMoments). The rates hold (Table~\ref{tab:repro}). The records
are read at \RecordsCommit.

\begin{table}[!htbp]
\centering\footnotesize
\ReproTable
\caption{The recorded baseline and kept variant beside the pinned app's texts on rebuilt evidence.}
\label{tab:repro}
\end{table}

\section{Limitations}

\begin{itemize}[nosep]
\item The grades check what can be checked: a true and concrete text can still fail to teach.
That question goes to the player on the Calibration page, and to the judge, whose agreement with
the player is itself only partly measured.
\item The judge is a model reading on the player's behalf. It is a good finder of what is
missing and a poorer judge of how much it matters, which is why the rounds use its pairwise
preferences and its notes, not its ratings.
\item \Games{} games of one player. The split keeps the claims honest about that player, not
about players in general.
\end{itemize}

\end{document}
